% ============================================================================
% Fichier  : partie4/P04_C12_artefact_vulnerabilite.tex
% Titre    : De l'Artefact à la Vulnérabilité
% Auteur   : MINKA MI NGUIDJOÏ Thierry Emmanuel
% Version  : 1.0 -- novembre 2026
% Corrigés : corriges/Corr_C12.tex
% ============================================================================
% Sources : ISO/IEC 27035 (enseignements tirés d'un incident, paraphrasée) ;
% NIST SP 800-61 r2 (activité post-incident) ; analyse des causes (arbre des
% causes, cinq pourquoi : méthodes générales, présentées sans citation).
% Faits nouveaux (bible) : export de configuration de la passerelle,
% micrologiciel 7.2.1 ; 4 requêtes du 02/03 entre 23 h 46 et 23 h 48 vers le
% chemin décrit comme indicateur par le bulletin de l'éditeur ; svc_backup
% membre du groupe des administrateurs du domaine ; aucune authentification
% à deux facteurs sur la passerelle ; clé de l'interface en clair dans
% api.ini ; 3 autres équipements de bordure non corrigés trouvés en
% septembre 2025.
% ============================================================================

\chapter{De l'Artefact à la Vulnérabilité}
\label{chap:artefact-vuln}
\chaptermark{De l'artefact à la vulnérabilité}

\epigraph{%
    « Trouver par où l'attaquant est entré, c'est la moitié de l'enquête.\\
    Comprendre pourquoi la porte était ouverte, c'est l'autre moitié. »%
}{Principe du cours}

\niveauChapitre{Intermédiaire à expert}{%
    Maîtrise des Chapitres~\ref{chap:chronologie} et~\ref{chap:gerer}
}

\begin{boxobjectifs}
\begin{itemize}
    \item Remonter d'un artefact au mécanisme qui l'a produit, et de ce
          mécanisme à une vulnérabilité identifiée.
    \item Confronter des traces aux indicateurs publiés par un éditeur.
    \item Examiner un équipement de bordure dont le système n'est pas
          accessible, et en dire les limites.
    \item Distinguer cause immédiate, causes contributives et cause
          profonde.
    \item Rechercher les variantes d'une faille dans le reste du parc.
    \item Transformer les enseignements d'une enquête en règles du programme
          de vulnérabilités.
\end{itemize}
\end{boxobjectifs}

\bigskip

% ============================================================================
\section*{Avant de commencer : une session sans mot de passe}
% ============================================================================

Yaoundé, 6 avril 2025. Trois semaines d'enquête ont établi la séquence de
l'attaque. Il reste un point que M.~ABOUBAKAR ne veut pas laisser à
l'inférence~: la première session du compte \texttt{svc\_backup}, le 2~mars,
s'est ouverte sans aucun échec d'authentification. Soit l'attaquant
connaissait le mot de passe, soit il n'en a pas eu besoin. La réponse
détermine la cause de l'incident, et elle pèse sur la question encore
ouverte de la participation d'un employé.

\begin{boxalert}
\textbf{Violation typique.} Désigner la faille la plus connue du moment
comme cause de l'incident, parce qu'elle touche le produit en cause, sans
confronter les traces à ses indicateurs.

\smallskip
\textbf{Règle~:} une vulnérabilité n'est la cause d'un incident que si les
traces portent sa signature, décrite par l'éditeur ou par un tiers
compétent, et si les autres explications sont écartées.
\end{boxalert}

\separateur

% ============================================================================
\section{Trois Explications pour une Session}
\label{sec:c12-explications}
% ============================================================================

\begin{tableaucours}{L{4.2cm}L{4.2cm}L{3.8cm}}{%
    \tableauHeader{Explication} &
    \tableauHeader{Traces attendues} &
    \tableauHeader{Traces observées}
}
E1~: mot de passe connu (fuite, employé, devinette) & Session directe, sans
requête inhabituelle avant & Requêtes inhabituelles avant la session \\
E2~: exploitation de VULN-2024-4471 & Requêtes vers le chemin décrit dans le
bulletin, puis session sans authentification & Quatre requêtes de ce type,
puis la session \\
E3~: vol de jeton d'une session légitime & Session légitime de
\texttt{svc\_backup} juste avant & Aucune session légitime depuis 9~jours \\
\end{tableaucours}
\vspace{-0.8em}
\begin{center}{\small\itshape Tableau~12.1. Trois explications de la session du 2~mars}\end{center}

Le bulletin publié par l'éditeur le 12~décembre 2024 décrit la faille, les
versions touchées (jusqu'à 7.2.4) et des \termecle{indicateurs de
compromission}~: la forme particulière des requêtes HTTP envoyées à la
passerelle lors de l'exploitation, et leur trace dans son journal. C'est la
signature à rechercher.

\begin{boxscenario}{Opération NYANGA. La signature du 2 mars}
\textbf{Constatations.}
\begin{itemize}
    \item L'export de configuration de la passerelle, réalisé le 14~mars
          (scellé NC-02), indique le micrologiciel 7.2.1, compris dans les
          versions touchées.
    \item Le journal de la passerelle enregistre, le 2~mars entre 23~h~46 et
          23~h~48 (heures corrigées), quatre requêtes depuis 203.0.113.18
          vers le chemin désigné par le bulletin comme indicateur
          d'exploitation, suivies de la session de \texttt{svc\_backup}.
    \item Aucune requête de ce type n'apparaît ailleurs dans les 60~jours
          du journal.
    \item Le mot de passe de \texttt{svc\_backup}, long de 24 caractères, n'a
          pas changé depuis 2022 et n'apparaît dans aucune fuite connue
          consultée.
\end{itemize}

\textbf{Analyse} (confiance élevée). Les traces portent la signature de
VULN-2024-4471 et précèdent immédiatement la session. E2 explique la session
sans authentification~; E1 ne l'explique pas sans supposer, en plus, que
l'attaquant connaissait un mot de passe qu'il n'a pas eu à saisir~; E3 est
écartée.
\end{boxscenario}

\begin{boxlecture}{Les trois explications au 6 avril}
\begin{tableaucours}{L{7.0cm}c}{%
    \tableauHeader{Proposition} &
    \tableauHeader{Cote}
}
E2~: la session du 2~mars résulte de l'exploitation de VULN-2024-4471 &
\OUI{} \\
E3~: la session résulte d'un vol de jeton d'une session légitime & \NON{} \\
E1~: l'attaquant connaissait le mot de passe de \texttt{svc\_backup} &
\NEK{} \\
Un employé a participé à l'attaque & \NEK{} \\
\end{tableaucours}

\medskip
\textbf{Lecture.} E1 reste \NEK{}~: qu'on n'ait pas eu besoin du mot de
passe ne prouve pas qu'on l'ignorait. Mais l'hypothèse perd son utilité~:
l'entrée s'explique sans elle. La participation d'un employé reste \NEK{},
et elle le restera tant que l'identité des auteurs l'est~; l'enquête ne
dispose d'aucune pièce qui la soutienne.
\end{boxlecture}

% ============================================================================
\section{Examiner un Équipement de Bordure}
\label{sec:c12-bordure}
% ============================================================================

Les passerelles d'accès, les pare-feu et les routeurs sont des appliances~:
leur système est fermé, souvent sans accès au disque ni à la mémoire.
L'enquêteur dispose de ce que l'équipement veut bien exporter.

\begin{tableaucours}{L{4.0cm}L{4.2cm}L{4.0cm}}{%
    \tableauHeader{Pièce} &
    \tableauHeader{Ce qu'elle apporte} &
    \tableauHeader{Limite}
}
Export de configuration & Version, comptes, règles, réglages & Instantané
au moment de l'export \\
Journaux exportés & Sessions, requêtes, erreurs & Rotation, niveau de détail
choisi \\
Outil d'intégrité de l'éditeur & Comparaison des fichiers système avec les
fichiers d'origine & Disponible pour certains produits seulement \\
Image du système, avec l'aide de l'éditeur & Examen complet & Rare, lent,
dépend de l'éditeur \\
\end{tableaucours}
\vspace{-0.8em}
\begin{center}{\small\itshape Tableau~12.2. Les pièces d'un équipement de bordure}\end{center}

Un attaquant qui contrôle une passerelle peut y avoir laissé un mécanisme de
retour qui survit à la mise à jour. Sans outil d'intégrité ni image du
système, la proposition « la passerelle ne contient plus de code de
l'attaquant » reste \NEK{}. À NYANGA, l'éditeur ne fournissait pas d'outil
d'intégrité pour la version 7.2~: la passerelle a été remplacée par un
équipement neuf, et l'ancienne conservée sous scellé.

% ============================================================================
\section{Cause Immédiate, Causes Contributives, Cause Profonde}
\label{sec:c12-causes}
% ============================================================================

\begin{boxdef}
\textbf{\termecle{Cause immédiate}}~: l'événement qui a directement permis
l'incident. \textbf{\termecle{Cause contributive}}~: condition qui a
aggravé l'incident ou en a facilité une étape. \textbf{\termecle{Cause
profonde}}~: défaut d'organisation qui a laissé subsister la cause
immédiate et les causes contributives~; la corriger empêche des incidents
semblables.
\end{boxdef}

La méthode des cinq pourquoi aide à descendre de la cause immédiate vers la
cause profonde~: on demande « pourquoi~? » à chaque réponse, jusqu'à
atteindre une décision ou une absence de décision de l'organisation.

\begin{boxscenario}{Opération NYANGA. Cinq pourquoi}
\begin{enumerate}
    \item Pourquoi l'attaquant est-il entré~? Parce que la passerelle
          présentait VULN-2024-4471.
    \item Pourquoi la faille était-elle présente~? Parce que le correctif du
          15~décembre n'avait pas été appliqué.
    \item Pourquoi n'avait-il pas été appliqué~? Parce que personne ne savait
          que la passerelle était concernée~: elle n'apparaissait dans aucun
          rapport de balayage.
    \item Pourquoi n'apparaissait-elle pas~? Parce qu'elle était absente de
          la base de gestion des configurations, qui fournissait la liste des
          cibles.
    \item Pourquoi était-elle absente~? Parce qu'aucune règle n'obligeait un
          prestataire à déclarer les équipements qu'il installe, ni personne
          à vérifier l'inventaire contre le réseau réel.
\end{enumerate}
\end{boxscenario}

\begin{boxlecture}{L'arbre des causes de NYANGA}
\begin{tableaucours}{L{6.2cm}L{2.4cm}c}{%
    \tableauHeader{Cause} &
    \tableauHeader{Type} &
    \tableauHeader{Établie}
}
VULN-2024-4471 non corrigée sur la passerelle & Immédiate & \OUI{} \\
\texttt{svc\_backup} membre des administrateurs du domaine & Contributive &
\OUI{} \\
Aucune authentification à deux facteurs sur la passerelle & Contributive &
\OUI{} \\
Identifiants d'administration présents en mémoire du contrôleur & Contributive
& \OUI{} \\
Clé de l'interface de paiement en clair dans \texttt{api.ini} & Contributive
& \OUI{} \\
Journaux de la passerelle conservés 60~jours & Contributive (enquête) &
\OUI{} \\
Inventaire non confronté au réseau réel, prestataires non tenus de
déclarer & Profonde & \OUI{} \\
\end{tableaucours}

\medskip
\textbf{Lecture.} Sept causes, toutes établies par des pièces. Corriger la
seule cause immédiate (appliquer le correctif) laisserait intactes les six
autres~: la prochaine faille d'un équipement non déclaré aurait le même
effet. La recommandation principale du rapport porte sur la cause profonde.
\end{boxlecture}

\begin{boxinfo}
\textbf{L'authentification à deux facteurs n'aurait pas suffi.}
VULN-2024-4471 contourne l'authentification~: un second facteur, vérifié par
le même mécanisme défaillant, aurait pu être contourné aussi. On la classe
en cause contributive parce qu'elle aurait gêné les sessions ultérieures sur
d'autres comptes, pas parce qu'elle aurait empêché l'entrée. Le rapport doit
le dire, pour ne pas vendre une mesure pour ce qu'elle n'est pas.
\end{boxinfo}

% ============================================================================
\section{Chercher les Variantes}
\label{sec:c12-variantes}
% ============================================================================

Une faille exploitée chez soi se cherche ensuite partout où elle peut
exister. C'est l'\termecle{analyse de variantes}~: la même vulnérabilité sur
d'autres équipements, la même faiblesse dans d'autres produits, la même
cause profonde ailleurs.

\begin{boxscenario}{Opération NYANGA. Septembre 2025}
\textbf{Constatations.} L'inventaire croisé de septembre 2025
(Chapitre~\ref{chap:gerer}) fait entrer quinze équipements dans la base.
Parmi eux, trois équipements de bordure (deux routeurs d'agence et un
point d'accès de télémaintenance) présentent des failles inscrites au
catalogue KEV, non corrigées depuis plus de six mois. Une recherche dans les
fichiers des serveurs trouve deux autres clés d'interface en clair, pour la
messagerie d'alerte et pour un service de change.

\textbf{Analyse} (confiance élevée). La cause profonde produisait ses effets
ailleurs~: trois autres portes ouvertes, deux autres clés exposées.
L'enquête sur un incident en a évité au moins un autre.
\end{boxscenario}

% ============================================================================
\section{Des Enseignements aux Règles}
\label{sec:c12-regles}
% ============================================================================

L'ISO/IEC~27035~\cite{iso27035} et le guide NIST SP~800-61~\cite{nist80061r2}
font des enseignements tirés d'un incident une étape à part entière. Ils ne
valent que s'ils deviennent des règles vérifiables.

\begin{tableaucours}{L{4.6cm}L{7.6cm}}{%
    \tableauHeader{Enseignement} &
    \tableauHeader{Règle du programme de NYANGA (septembre 2025)}
}
Un équipement non déclaré n'est pas analysé & Découverte réseau mensuelle de
toutes les plages, écarts avec la base traités sous 15~jours \\
Les équipements de bordure sont la porte d'entrée & Toute faille inscrite au
catalogue KEV sur un équipement exposé est classée P1, délai 72~heures \\
Un prestataire installe sans déclarer & Clause contractuelle de déclaration
et d'accès au balayage \\
Les journaux de la passerelle ont failli disparaître & Conservation de dix
ans des données de connexion (loi n°~2010/012, art.~29) \\
Une clé en clair a permis les virements & Coffre de secrets pour toutes les
clés d'interface, recherche trimestrielle de secrets en clair \\
\end{tableaucours}
\vspace{-0.8em}
\begin{center}{\small\itshape Tableau~12.3. Des enseignements aux règles}\end{center}

% ============================================================================
\section{Ce que la Remontée demande à chaque Métier}
\label{sec:c12-metiers}
% ============================================================================

\begin{boxmetier}{Analyste en réponse à incident}
\textbf{Remonter à la vulnérabilité}
\begin{enumerate}
    \item Lister toutes les explications du premier accès.
    \item Pour chacune, écrire les traces attendues.
    \item Confronter les traces aux indicateurs publiés par l'éditeur.
    \item Coter chaque explication, et écrire ce qui reste \NEK{}.
    \item Construire l'arbre des causes jusqu'à la cause profonde.
    \item Lancer la recherche de variantes dans tout le parc.
\end{enumerate}
\end{boxmetier}

\begin{boxmetier}{Analyste en gestion des vulnérabilités}
Recevoir chaque rapport d'incident comme une mesure de son programme~: la
faille exploitée était-elle connue, classée, dans le délai~? Chaque réponse
« non » appelle une règle nouvelle.
\end{boxmetier}

\begin{boxmetier}{Application OPJ/Expert Judiciaire}
L'expert établit le mécanisme d'entrée parce qu'il répond à l'article~69
de la loi n°~2010/012 (violation des mesures de sécurité). Il ne juge pas
la diligence de la victime, sauf si sa mission le lui demande.
\end{boxmetier}

\begin{boxmetier}{RSSI}
Présenter à la direction l'arbre des causes complet, et faire décider la
correction de la cause profonde. Un correctif appliqué en une heure ne
clôt pas l'incident~; une règle d'inventaire tenue chaque mois le clôt.
\end{boxmetier}

\separateur

\begin{boxresume}
\textbf{Ce qu'il faut retenir de ce chapitre~:}
\begin{itemize}
    \item Une vulnérabilité est la cause d'un incident quand les traces
          portent sa signature et que les autres explications sont
          écartées.
    \item Un équipement de bordure s'examine par ses exports~; ce qu'on ne
          peut pas lire reste \NEK{}.
    \item On distingue cause immédiate, causes contributives et cause
          profonde~; la recommandation principale vise la cause profonde.
    \item L'analyse de variantes cherche la même faille et la même cause
          ailleurs.
    \item Un enseignement ne vaut que transformé en règle vérifiable.
\end{itemize}
\end{boxresume}

\begin{boxinfo}
\textbf{Points de contrôle.} Vous avez compris ce chapitre si vous savez~:
\begin{itemize}
    \item confronter des traces aux indicateurs d'un bulletin~;
    \item construire un arbre des causes à trois niveaux~;
    \item rédiger une règle de programme à partir d'un enseignement.
\end{itemize}
\end{boxinfo}

% ============================================================================
\section*{Exercices}
% ============================================================================

\begin{boxexo}[title={\ding{45}\quad Exercice~12.1. Immédiate, contributive, profonde \niveaubase}]
Classez~: (a) le correctif non appliqué~; (b) l'absence de règle de
déclaration des équipements~; (c) la clé en clair dans \texttt{api.ini}~;
(d) le compte \texttt{svc\_backup} trop privilégié~; (e) la durée de
conservation des journaux.
\end{boxexo}

\begin{boxexo}[title={\ding{45}\quad Exercice~12.2. Une autre signature \niveaumoyen}]
Supposez que le journal de la passerelle n'ait montré aucune requête vers le
chemin du bulletin, mais trois échecs d'authentification de
\texttt{svc\_backup} depuis 203.0.113.18, puis une session réussie.
\begin{enumerate}[(a)]
    \item Comment cotez-vous E1 et E2~?
    \item Quelle cause immédiate retenez-vous alors, et quelles causes
          contributives changent~?
    \item Que devient la question de la participation d'un employé~?
\end{enumerate}
\end{boxexo}

\begin{boxexo}[title={\ding{45}\quad Exercice~12.3. Les cinq pourquoi de SOLITRANS \niveaumoyen}]
Le terminal de pesage de Bonabéri a cessé de répondre pendant un balayage
(Chapitre~\ref{chap:identifier}). Appliquez la méthode des cinq pourquoi et
identifiez la cause profonde. Proposez la règle qui en découle.
\end{boxexo}

\begin{boxexo}[title={\ding{45}\quad Exercice~12.4. La passerelle sous scellé \niveauexpert}]
\begin{enumerate}[(a)]
    \item Pourquoi a-t-on remplacé la passerelle au lieu de simplement la
          mettre à jour~?
    \item Quelles questions l'ancienne passerelle, conservée sous scellé,
          pourrait-elle encore trancher, et avec l'aide de qui~?
    \item Rédigez le paragraphe du rapport d'expertise qui établit la
          violation des mesures de sécurité au sens de l'article~69, en
          séparant constatation et analyse.
\end{enumerate}
\end{boxexo}

\bigskip

\begin{boxinfo}
\textbf{Pour aller plus loin.}
\begin{itemize}
    \item L'ISO/IEC~27035~\cite{iso27035}, sur la gestion des incidents et
          leurs enseignements.
    \item Le NIST SP~800-61~\cite{nist80061r2}, section sur l'activité
          post-incident.
\end{itemize}
\end{boxinfo}

\begin{boxinfo}
\textbf{Transition.} L'attaquant de NYANGA n'a presque rien effacé. Un
attaquant plus soigneux maquille ses dates, vide les journaux, ou travaille
dans des systèmes que l'enquêteur ne touche pas~: Chapitre~\ref{chap:antiforensique}.
\end{boxinfo}
